It's surprising that, despite all the expressivity of the univariate TNBBeta distribution,
there is little difference between the latent spaces in both hyperspherical VAEs tested above.
It would be worth understanding whether this is a theoretical necessity, or if this is simply
a domain-specific quirk that may not hold in other datasets. I tried fixing some parameters
to see if I could force ring-like posteriors. Fixing $\boldsymbol{\mu}$ only resulted in
situations in which $p \to 0$ and $q \to 1$. When I fixed $\varepsilon$, $q$ replaced $\varepsilon$
as a concentratio parameter in the posterior. Essentially, I have a sense that the S-VAE and
TNBBeta-VAE are producing highly similar posteriors in low dimensions, and I would like to understand
why.

Additionally, the SphericalTNBBeta VAE presents an alternative to the S-VAE that has faster sampling
and is more stable in high dimensions. This is the same conclusion as \citep{decao2020powersphericaldistribution}'s
power spherical distribution. In this light, it is worth understanding how these two hyperspherical
approaches to VAEs differ.